% FILE: CSE-23_TermFinal.tex

\documentclass[11pt, a4paper]{article}

% --- UNIVERSAL PREAMBLE BLOCK ---
\usepackage[a4paper, top=2.5cm, bottom=2.5cm, left=2cm, right=2cm]{geometry}
\usepackage{fontspec}
\usepackage[english, bidi=basic, provide=*]{babel}

% Set default/Latin font to Sans Serif in the main (rm) slot
\babelfont{rm}{Noto Sans}

% --- EXTRA PACKAGES FOR MATHEMATICS AND FORMATTING ---
\usepackage{amsmath}
\usepackage{amssymb}
\usepackage{enumitem}

\begin{document}

%%%%%%%%%%%%%%%%%%%%%%%%%%%%%%%%%%%%%%%%%%%%%%%%%%
% QUESTION_START
%%%%%%%%%%%%%%%%%%%%%%%%%%%%%%%%%%%%%%%%%%%%%%%%%%

% id=cse23-final-secb-q5a
% discipline=CSE
% year=2023
% exam_type=Term Final
% section=B
% ct_number=UNKNOWN
% question_number=5a
% topic=Definite Integration
% subtopic=General
% difficulty=Easy
% length=Short
% question_type=Repeated PYQ Pattern
% frequency=1
% appearances=
% OCR_STATUS=VERIFIED

\section*{Term Final - Q5(a)}

Question:
What is anti-derivative? Find a function $f(x)$ such that $f''(x) = x + \cos x$ where $f(0) = 1$ and $f'(0) = 2$. [Marks: 04]

Solution:
Step 1: Define anti-derivative.
An antiderivative of a function $f(x)$ on an interval $I$ is a differentiable function $F(x)$ whose derivative is equal to the original function $f(x)$ for all $x \in I$. That is:
\[
F'(x) = f(x) \quad \text{for all } x \in I
\]

Step 2: Find $f'(x)$ by integrating $f''(x)$.
Given the second derivative:
\[
f''(x) = x + \cos x
\]
Integrating both sides once with respect to $x$ gives:
\[
f'(x) = \int (x + \cos x) \, dx = \frac{x^2}{2} + \sin x + C_1
\]
where $C_1$ is an arbitrary constant of integration.

Step 3: Use the initial condition $f'(0) = 2$ to find $C_1$.
\[
f'(0) = \frac{0^2}{2} + \sin(0) + C_1 = 2 \implies C_1 = 2
\]
Thus, the first derivative is:
\[
f'(x) = \frac{x^2}{2} + \sin x + 2
\]

Step 4: Find $f(x)$ by integrating $f'(x)$.
Integrating $f'(x)$ with respect to $x$:
\[
f(x) = \int \left( \frac{x^2}{2} + \sin x + 2 \right) \, dx = \frac{x^3}{6} - \cos x + 2x + C_2
\]
where $C_2$ is another constant of integration.

Step 5: Use the initial condition $f(0) = 1$ to find $C_2$.
\[
f(0) = \frac{0^3}{6} - \cos(0) + 2(0) + C_2 = 1 \implies -1 + C_2 = 1 \implies C_2 = 2
\]
Therefore, the function is:
\[
f(x) = \frac{x^3}{6} - \cos x + 2x + 2
\]

Final Answer:
\[
\boxed{f(x) = \frac{x^3}{6} - \cos x + 2x + 2}
\]

%%%%%%%%%%%%%%%%%%%%%%%%%%%%%%%%%%%%%%%%%%%%%%%%%%
% QUESTION_END
%%%%%%%%%%%%%%%%%%%%%%%%%%%%%%%%%%%%%%%%%%%%%%%%%%


%%%%%%%%%%%%%%%%%%%%%%%%%%%%%%%%%%%%%%%%%%%%%%%%%%
% QUESTION_START
%%%%%%%%%%%%%%%%%%%%%%%%%%%%%%%%%%%%%%%%%%%%%%%%%%

% id=cse23-final-secb-q5b
% discipline=CSE
% year=2023
% exam_type=Term Final
% section=B
% ct_number=UNKNOWN
% question_number=5b
% topic=Indefinite Integration
% subtopic=Reduction Formula
% difficulty=Medium
% length=Long
% question_type=Repeated PYQ Pattern
% frequency=1
% appearances=
% OCR_STATUS=VERIFIED

\section*{Term Final - Q5(b)}

Question:
Obtain the reduction formula for $\int \cos^n x \, dx$ and hence deduce the integral $\int_{0}^{\frac{\pi}{2}} \cos^9 x \, dx$. [Marks: 06]

Solution:
Step 1: Obtain the reduction formula.
Let $J_n = \int \cos^n x \, dx$. Rewrite the integrand:
\[
J_n = \int \cos^{n-1} x \cdot \cos x \, dx
\]
Apply integration by parts: $\int u \, dv = uv - \int v \, du$.
Let $u = \cos^{n-1} x \implies du = -(n-1)\cos^{n-2} x \sin x \, dx$.
Let $dv = \cos x \, dx \implies v = \sin x$.
Applying the integration by parts formula:
\[
J_n = \cos^{n-1} x \sin x - \int \sin x \cdot \left[ -(n-1)\cos^{n-2} x \sin x \right] \, dx
\]
\[
J_n = \cos^{n-1} x \sin x + (n-1) \int \cos^{n-2} x \sin^2 x \, dx
\]
Substitute $\sin^2 x = 1 - \cos^2 x$:
\[
J_n = \cos^{n-1} x \sin x + (n-1) \int \cos^{n-2} x (1 - \cos^2 x) \, dx
\]
\[
J_n = \cos^{n-1} x \sin x + (n-1) \int \cos^{n-2} x \, dx - (n-1) \int \cos^n x \, dx
\]
Using our notation $J_n$:
\[
J_n = \cos^{n-1} x \sin x + (n-1) J_{n-2} - (n-1) J_n
\]
Rearranging and grouping $J_n$ terms:
\[
J_n + (n-1) J_n = \cos^{n-1} x \sin x + (n-1) J_{n-2} \implies n J_n = \cos^{n-1} x \sin x + (n-1) J_{n-2}
\]
Thus, the reduction formula is:
\[
J_n = \frac{\cos^{n-1} x \sin x}{n} + \frac{n-1}{n} J_{n-2}
\]

Step 2: Adapt the reduction formula to the definite integral $I_n = \int_{0}^{\frac{\pi}{2}} \cos^n x \, dx$.
\[
I_n = \left[ \frac{\cos^{n-1} x \sin x}{n} \right]_{0}^{\frac{\pi}{2}} + \frac{n-1}{n} I_{n-2}
\]
Since $\cos(\frac{\pi}{2}) = 0$ (for $n \ge 2$) and $\sin(0) = 0$, the boundary term evaluates to $0$:
\[
I_n = \frac{n-1}{n} I_{n-2}
\]

Step 3: Evaluate $I_9 = \int_{0}^{\frac{\pi}{2}} \cos^9 x \, dx$.
Applying the recurrence relation repeatedly:
\[
I_9 = \frac{8}{9} I_7 = \frac{8}{9} \left( \frac{6}{7} I_5 \right) = \frac{8}{9} \cdot \frac{6}{7} \cdot \frac{4}{5} I_3 = \frac{8}{9} \cdot \frac{6}{7} \cdot \frac{4}{5} \cdot \frac{2}{3} I_1
\]
Compute the base case $I_1$:
\[
I_1 = \int_{0}^{\frac{\pi}{2}} \cos x \, dx = \left[ \sin x \right]_{0}^{\frac{\pi}{2}} = \sin\left(\frac{\pi}{2}\right) - \sin(0) = 1
\]
Therefore,
\[
I_9 = \frac{8 \cdot 6 \cdot 4 \cdot 2}{9 \cdot 7 \cdot 5 \cdot 3} \cdot 1 = \frac{384}{945} = \frac{128}{315}
\]

Final Answer:
\[
\boxed{\frac{128}{315}}
\]

%%%%%%%%%%%%%%%%%%%%%%%%%%%%%%%%%%%%%%%%%%%%%%%%%%
% QUESTION_END
%%%%%%%%%%%%%%%%%%%%%%%%%%%%%%%%%%%%%%%%%%%%%%%%%%


%%%%%%%%%%%%%%%%%%%%%%%%%%%%%%%%%%%%%%%%%%%%%%%%%%
% QUESTION_START
%%%%%%%%%%%%%%%%%%%%%%%%%%%%%%%%%%%%%%%%%%%%%%%%%%

% id=cse23-final-secb-q6a
% discipline=CSE
% year=2023
% exam_type=Term Final
% section=B
% ct_number=UNKNOWN
% question_number=6a
% topic=Definite Integration
% subtopic=General
% difficulty=Medium
% length=Long
% question_type=Repeated PYQ Pattern
% frequency=1
% appearances=
% OCR_STATUS=VERIFIED

\section*{Term Final - Q6(a)}

Question:
State and prove the fundamental theorem of calculus. [Marks: 04]

Solution:
Step 1: State the Fundamental Theorem of Calculus (FTC).
The Fundamental Theorem of Calculus consists of two parts:

**Part 1:** If $f$ is continuous on the closed interval $[a, b]$, then the function $g$ defined by:
\[
g(x) = \int_{a}^{x} f(t) \, dt \quad a \le x \le b
\]
is continuous on $[a, b]$, differentiable on $(a, b)$, and its derivative is:
\[
g'(x) = f(x)
\]

**Part 2:** If $f$ is continuous on $[a, b]$, and $F$ is any antiderivative of $f$ (so that $F'(x) = f(x)$), then:
\[
\int_{a}^{b} f(x) \, dx = F(b) - F(a)
\]

Step 2: Prove Part 1 of the theorem.
By the definition of the derivative:
\[
g'(x) = \lim_{h \to 0} \frac{g(x+h) - g(x)}{h}
\]
Substituting the definition of $g(x)$:
\[
g(x+h) - g(x) = \int_{a}^{x+h} f(t) \, dt - \int_{a}^{x} f(t) \, dt = \int_{x}^{x+h} f(t) \, dt
\]
For $h \ne 0$, we have:
\[
\frac{g(x+h) - g(x)}{h} = \frac{1}{h} \int_{x}^{x+h} f(t) \, dt
\]
By the Mean Value Theorem for Integrals, if $f$ is continuous on $[x, x+h]$, there exists a number $c$ in this interval such that:
\[
\int_{x}^{x+h} f(t) \, dt = f(c) h
\]
Thus:
\[
\frac{g(x+h) - g(x)}{h} = \frac{1}{h} (f(c) h) = f(c)
\]
Taking the limit as $h \to 0$, the interval $[x, x+h]$ shrinks to the single point $x$, forcing $c \to x$. Since $f$ is continuous, we get:
\[
g'(x) = \lim_{h \to 0} f(c) = f(x)
\]
This completes the proof of Part 1.

Step 3: Prove Part 2 of the theorem.
Let $g(x) = \int_{a}^{x} f(t) \, dt$. From Part 1, we know that $g'(x) = f(x)$.
If $F$ is any other antiderivative of $f$, then $F'(x) = f(x) = g'(x)$.
Thus, $g(x)$ and $F(x)$ differ by a constant $C$:
\[
g(x) = F(x) + C \quad \text{for } a \le x \le b
\]
To find $C$, evaluate this equation at $x = a$:
\[
g(a) = F(a) + C \implies \int_{a}^{a} f(t) \, dt = F(a) + C \implies 0 = F(a) + C \implies C = -F(a)
\]
Substitute $C$ back into the relation:
\[
g(x) = F(x) - F(a)
\]
Now evaluate at $x = b$:
\[
g(b) = F(b) - F(a) \implies \int_{a}^{b} f(x) \, dx = F(b) - F(a)
\]
This completes the proof of Part 2.

Final Answer:
\[
\boxed{\text{FTC Part 1: } \frac{d}{dx}\int_a^x f(t)dt = f(x); \text{ Part 2: } \int_a^b f(x)dx = F(b)-F(a). \text{ Proved using MVT for integrals and antiderivatives.}}
\]

%%%%%%%%%%%%%%%%%%%%%%%%%%%%%%%%%%%%%%%%%%%%%%%%%%
% QUESTION_END
%%%%%%%%%%%%%%%%%%%%%%%%%%%%%%%%%%%%%%%%%%%%%%%%%%


%%%%%%%%%%%%%%%%%%%%%%%%%%%%%%%%%%%%%%%%%%%%%%%%%%
% QUESTION_START
%%%%%%%%%%%%%%%%%%%%%%%%%%%%%%%%%%%%%%%%%%%%%%%%%%

% id=cse23-final-secb-q6b
% discipline=CSE
% year=2023
% exam_type=Term Final
% section=B
% ct_number=UNKNOWN
% question_number=6b
% topic=Definite Integration
% subtopic=General
% difficulty=Medium
% length=Medium
% question_type=Repeated PYQ Pattern
% frequency=1
% appearances=
% OCR_STATUS=VERIFIED

\section*{Term Final - Q6(b)}

Question:
Evaluate the integral $\int_{2}^{3} x^2 \, dx$ as the limit of a sum. [Marks: 03]

Solution:
Step 1: Set up the limit of a sum definition of a definite integral.
The definition of the definite integral as the limit of a Riemann sum is:
\[
\int_{a}^{b} f(x) \, dx = \lim_{n \to \infty} \sum_{i=1}^{n} f(a + i \Delta x) \Delta x
\]
where $a = 2$, $b = 3$, $f(x) = x^2$, and:
\[
\Delta x = \frac{b-a}{n} = \frac{3-2}{n} = \frac{1}{n}
\]
The evaluation points are:
\[
x_i = a + i \Delta x = 2 + \frac{i}{n}
\]

Step 2: Express the function at the evaluation points.
\[
f(x_i) = \left( 2 + \frac{i}{n} \right)^2 = 4 + \frac{4i}{n} + \frac{i^2}{n^2}
\]

Step 3: Write down the summation expression.
\[
S_n = \sum_{i=1}^{n} \left( 4 + \frac{4i}{n} + \frac{i^2}{n^2} \right) \frac{1}{n} = \frac{4}{n} \sum_{i=1}^{n} 1 + \frac{4}{n^2} \sum_{i=1}^{n} i + \frac{1}{n^3} \sum_{i=1}^{n} i^2
\]
Using the standard summation formulas:
\[
\sum_{i=1}^{n} 1 = n, \quad \sum_{i=1}^{n} i = \frac{n(n+1)}{2}, \quad \sum_{i=1}^{n} i^2 = \frac{n(n+1)(2n+1)}{6}
\]
Substitute these into the expression:
\[
S_n = \frac{4}{n}(n) + \frac{4}{n^2} \left[ \frac{n(n+1)}{2} \right] + \frac{1}{n^3} \left[ \frac{n(n+1)(2n+1)}{6} \right]
\]
\[
S_n = 4 + 2 \left( 1 + \frac{1}{n} \right) + \frac{1}{6} \left( 1 + \frac{1}{n} \right) \left( 2 + \frac{1}{n} \right)
\]

Step 4: Take the limit of $S_n$ as $n \to \infty$.
\[
\lim_{n \to \infty} S_n = 4 + 2(1) + \frac{1}{6}(1)(2) = 4 + 2 + \frac{1}{3} = 6 + \frac{1}{3} = \frac{19}{3}
\]

Final Answer:
\[
\boxed{\frac{19}{3}}
\]

%%%%%%%%%%%%%%%%%%%%%%%%%%%%%%%%%%%%%%%%%%%%%%%%%%
% QUESTION_END
%%%%%%%%%%%%%%%%%%%%%%%%%%%%%%%%%%%%%%%%%%%%%%%%%%


%%%%%%%%%%%%%%%%%%%%%%%%%%%%%%%%%%%%%%%%%%%%%%%%%%
% QUESTION_START
%%%%%%%%%%%%%%%%%%%%%%%%%%%%%%%%%%%%%%%%%%%%%%%%%%

% id=cse23-final-secb-q6c
% discipline=CSE
% year=2023
% exam_type=Term Final
% section=B
% ct_number=UNKNOWN
% question_number=6c
% topic=Definite Integration
% subtopic=General
% difficulty=Medium
% length=Medium
% question_type=Repeated PYQ Pattern
% frequency=1
% appearances=
% OCR_STATUS=VERIFIED

\section*{Term Final - Q6(c)}

Question:
Prove that $\int_{0}^{\pi} \frac{x \, dx}{1 + \sin x} = \pi$. [Marks: 03]

Solution:
Step 1: Set up the integral and apply the definite integral property.
Let:
\[
I = \int_{0}^{\pi} \frac{x \, dx}{1 + \sin x} \quad \text{--- (1)}
\]
Using the property $\int_{a}^{b} f(x) \, dx = \int_{a}^{b} f(a+b-x) \, dx$:
\[
I = \int_{0}^{\pi} \frac{(\pi - x) \, dx}{1 + \sin(\pi - x)} = \int_{0}^{\pi} \frac{(\pi - x) \, dx}{1 + \sin x} \quad \text{--- (2)}
\]

Step 2: Add equations (1) and (2).
\[
2I = \int_{0}^{\pi} \frac{x + \pi - x}{1 + \sin x} \, dx = \pi \int_{0}^{\pi} \frac{dx}{1 + \sin x}
\]

Step 3: Evaluate the integral $J = \int_{0}^{\pi} \frac{dx}{1 + \sin x}$.
To handle this strictly and avoid the discontinuity of standard trigonometric transformations at $\frac{\pi}{2}$, we use Weierstrass substitution ($t = \tan \frac{x}{2}$):
\[
dx = \frac{2 \, dt}{1 + t^2}, \quad \sin x = \frac{2t}{1 + t^2}
\]
Determine the new boundaries of integration:
\begin{itemize}
    \item When $x = 0 \implies t = \tan(0) = 0$.
    \item When $x \to \pi \implies t \to \infty$.
\end{itemize}
Substitute these into the integral:
\[
J = \int_{0}^{\infty} \frac{\frac{2 \, dt}{1 + t^2}}{1 + \frac{2t}{1 + t^2}} = \int_{0}^{\infty} \frac{2 \, dt}{t^2 + 2t + 1} = \int_{0}^{\infty} \frac{2 \, dt}{(t + 1)^2}
\]
Evaluate the improper integral:
\[
J = \left[ -\frac{2}{t + 1} \right]_{0}^{\infty} = 0 - (-2) = 2
\]

Step 4: Solve for $I$.
Since $2I = \pi J$:
\[
2I = \pi (2) \implies I = \pi \quad \text{(Proved)}
\]

Final Answer:
\[
\boxed{\pi}
\]

%%%%%%%%%%%%%%%%%%%%%%%%%%%%%%%%%%%%%%%%%%%%%%%%%%
% QUESTION_END
%%%%%%%%%%%%%%%%%%%%%%%%%%%%%%%%%%%%%%%%%%%%%%%%%%


%%%%%%%%%%%%%%%%%%%%%%%%%%%%%%%%%%%%%%%%%%%%%%%%%%
% QUESTION_START
%%%%%%%%%%%%%%%%%%%%%%%%%%%%%%%%%%%%%%%%%%%%%%%%%%

% id=cse23-final-secb-q7a
% discipline=CSE
% year=2023
% exam_type=Term Final
% section=B
% ct_number=UNKNOWN
% question_number=7a
% topic=Definite Integration
% subtopic=General
% difficulty=Medium
% length=Medium
% question_type=Repeated PYQ Pattern
% frequency=1
% appearances=
% OCR_STATUS=VERIFIED

\section*{Term Final - Q7(a)}

Question:
Define Beta and Gamma functions. Prove that $\beta(m, n) = \frac{\Gamma(m)\Gamma(n)}{\Gamma(m+n)}$. [Marks: 04]

Solution:
Step 1: Define Beta and Gamma functions.
1. **Gamma Function ($\Gamma(n)$):** For $n > 0$, the Gamma function is defined as:
\[
\Gamma(n) = \int_{0}^{\infty} x^{n-1} e^{-x} \, dx
\]
2. **Beta Function ($\beta(m, n)$):** For $m, n > 0$, the Beta function is defined as:
\[
\beta(m, n) = \int_{0}^{1} x^{m-1} (1-x)^{n-1} \, dx
\]

Step 2: Express $\Gamma(m)$ and $\Gamma(n)$ in a substituted double-integral format.
Using the substitution $x = u^2 \implies dx = 2u \, du$:
\[
\Gamma(m) = 2 \int_{0}^{\infty} u^{2m-1} e^{-u^2} \, du \quad \text{--- (1)}
\]
Similarly, for $\Gamma(n)$ using $y = v^2$:
\[
\Gamma(n) = 2 \int_{0}^{\infty} v^{2n-1} e^{-v^2} \, dv \quad \text{--- (2)}
\]

Step 3: Multiply equations (1) and (2) and switch to polar coordinates.
\[
\Gamma(m)\Gamma(n) = 4 \int_{0}^{\infty} \int_{0}^{\infty} u^{2m-1} v^{2n-1} e^{-(u^2 + v^2)} \, du \, dv
\]
Let $u = r \cos \theta$, $v = r \sin \theta$. The area element becomes $du \, dv = r \, dr \, d\theta$, and boundaries are $r \in [0, \infty)$, $\theta \in [0, \frac{\pi}{2}]$:
\[
\Gamma(m)\Gamma(n) = 4 \int_{0}^{\frac{\pi}{2}} \int_{0}^{\infty} (r \cos \theta)^{2m-1} (r \sin \theta)^{2n-1} e^{-r^2} r \, dr \, d\theta
\]
\[
= 2 \left( 2 \int_{0}^{\frac{\pi}{2}} \cos^{2m-1}\theta \sin^{2n-1}\theta \, d\theta \right) \left( \int_{0}^{\infty} r^{2(m+n)-1} e^{-r^2} \, dr \right)
\]

Step 4: Connect the expressions back to the functions.
Recall the trigonometric form of the Beta function:
\[
\beta(m, n) = 2 \int_{0}^{\frac{\pi}{2}} \sin^{2n-1}\theta \cos^{2m-1}\theta \, d\theta
\]
Also, let $t = r^2 \implies dt = 2r \, dr$ in the radial integral:
\[
\int_{0}^{\infty} r^{2(m+n)-1} e^{-r^2} \, dr = \frac{1}{2} \int_{0}^{\infty} t^{(m+n)-1} e^{-t} \, dt = \frac{1}{2} \Gamma(m+n)
\]
Substitute these into our product formula:
\[
\Gamma(m)\Gamma(n) = 2 \left[ \beta(m, n) \right] \left[ \frac{1}{2} \Gamma(m+n) \right] = \beta(m, n) \Gamma(m+n)
\]
Thus:
\[
\beta(m, n) = \frac{\Gamma(m)\Gamma(n)}{\Gamma(m+n)} \quad \text{(Proved)}
\]

Final Answer:
\[
\boxed{\beta(m, n) = \frac{\Gamma(m)\Gamma(n)}{\Gamma(m+n)}}
\]

%%%%%%%%%%%%%%%%%%%%%%%%%%%%%%%%%%%%%%%%%%%%%%%%%%
% QUESTION_END
%%%%%%%%%%%%%%%%%%%%%%%%%%%%%%%%%%%%%%%%%%%%%%%%%%


%%%%%%%%%%%%%%%%%%%%%%%%%%%%%%%%%%%%%%%%%%%%%%%%%%
% QUESTION_START
%%%%%%%%%%%%%%%%%%%%%%%%%%%%%%%%%%%%%%%%%%%%%%%%%%

% id=cse23-final-secb-q7b
% discipline=CSE
% year=2023
% exam_type=Term Final
% section=B
% ct_number=UNKNOWN
% question_number=7b
% topic=Definite Integration
% subtopic=General
% difficulty=Medium
% length=Medium
% question_type=Repeated PYQ Pattern
% frequency=1
% appearances=
% OCR_STATUS=VERIFIED

\section*{Term Final - Q7(b)}

Question:
Evaluate $\int_{0}^{1} x^4 (1 - \sqrt{x})^5 \, dx$ by using gamma-beta function. [Marks: 03]

Solution:
Step 1: Simplify the integral using substitution.
Let $u = \sqrt{x} \implies x = u^2$.
Differentiating:
\[
dx = 2u \, du
\]
The integration boundaries remain $0$ and $1$:
\begin{itemize}
    \item When $x = 0 \implies u = 0$.
    \item When $x = 1 \implies u = 1$.
\end{itemize}

Step 2: Substitute into the integral.
\[
I = \int_{0}^{1} x^4 (1 - \sqrt{x})^5 \, dx = \int_{0}^{1} (u^2)^4 (1 - u)^5 (2u \, du)
\]
\[
= 2 \int_{0}^{1} u^8 (1 - u)^5 u \, du = 2 \int_{0}^{1} u^9 (1 - u)^5 \, du
\]

Step 3: Relate to the Beta function definition.
The definition of the Beta function is $\beta(p, q) = \int_{0}^{1} u^{p-1} (1-u)^{q-1} \, du$.
Here:
\[
p - 1 = 9 \implies p = 10
\]
\[
q - 1 = 5 \implies q = 6
\]
So,
\[
I = 2 \beta(10, 6)
\]

Step 4: Evaluate using the Beta-Gamma relation.
\[
I = 2 \frac{\Gamma(10)\Gamma(6)}{\Gamma(16)} = 2 \frac{9! \cdot 5!}{15!} = 2 \frac{9! \cdot 120}{15 \cdot 14 \cdot 13 \cdot 12 \cdot 11 \cdot 10 \cdot 9!}
\]
Simplifying the fractions:
\[
I = 2 \frac{120}{15 \cdot 14 \cdot 13 \cdot 12 \cdot 11 \cdot 10} = \frac{240}{3603600} = \frac{1}{15015}
\]

Final Answer:
\[
\boxed{\frac{1}{15015}}
\]

%%%%%%%%%%%%%%%%%%%%%%%%%%%%%%%%%%%%%%%%%%%%%%%%%%
% QUESTION_END
%%%%%%%%%%%%%%%%%%%%%%%%%%%%%%%%%%%%%%%%%%%%%%%%%%


%%%%%%%%%%%%%%%%%%%%%%%%%%%%%%%%%%%%%%%%%%%%%%%%%%
% QUESTION_START
%%%%%%%%%%%%%%%%%%%%%%%%%%%%%%%%%%%%%%%%%%%%%%%%%%

% id=cse23-final-secb-q7c
% discipline=CSE
% year=2023
% exam_type=Term Final
% section=B
% ct_number=UNKNOWN
% question_number=7c
% topic=Definite Integration
% subtopic=General
% difficulty=Easy
% length=Short
% question_type=Repeated PYQ Pattern
% frequency=1
% appearances=
% OCR_STATUS=VERIFIED

\section*{Term Final - Q7(c)}

Question:
Determine whether the integral $\int_{1}^{\infty} \frac{dx}{x^2 + 1}$ converges or diverges. [Marks: 03]

Solution:
Step 1: Set up the limit for the improper integral.
By the definition of an improper integral with an infinite upper bound:
\[
\int_{1}^{\infty} \frac{dx}{x^2 + 1} = \lim_{b \to \infty} \int_{1}^{b} \frac{dx}{x^2 + 1}
\]

Step 2: Integrate using the standard inverse tangent antiderivative.
\[
\lim_{b \to \infty} \left[ \tan^{-1} x \right]_{1}^{b} = \lim_{b \to \infty} \left( \tan^{-1} b - \tan^{-1} 1 \right)
\]

Step 3: Evaluate the limits.
We know that:
\[
\lim_{b \to \infty} \tan^{-1} b = \frac{\pi}{2}
\]
\[
\tan^{-1} 1 = \frac{\pi}{4}
\]
Thus:
\[
\int_{1}^{\infty} \frac{dx}{x^2 + 1} = \frac{\pi}{2} - \frac{\pi}{4} = \frac{\pi}{4}
\]
Since the limit yields a finite value, the improper integral converges.

Final Answer:
\[
\boxed{\text{The improper integral converges to } \frac{\pi}{4}}
\]

%%%%%%%%%%%%%%%%%%%%%%%%%%%%%%%%%%%%%%%%%%%%%%%%%%
% QUESTION_END
%%%%%%%%%%%%%%%%%%%%%%%%%%%%%%%%%%%%%%%%%%%%%%%%%%


%%%%%%%%%%%%%%%%%%%%%%%%%%%%%%%%%%%%%%%%%%%%%%%%%%
% QUESTION_START
%%%%%%%%%%%%%%%%%%%%%%%%%%%%%%%%%%%%%%%%%%%%%%%%%%

% id=cse23-final-secb-q8a
% discipline=CSE
% year=2023
% exam_type=Term Final
% section=B
% ct_number=UNKNOWN
% question_number=8a
% topic=Definite Integration
% subtopic=General
% difficulty=Medium
% length=Medium
% question_type=Repeated PYQ Pattern
% frequency=1
% appearances=
% OCR_STATUS=VERIFIED

\section*{Term Final - Q8(a)}

Question:
Sketch the region enclosed by the curves $y = x^2$ and $y = \sqrt{x}$, $x = \frac{1}{4}$, $x = 1$ and find its area. [Marks: 05]

Solution:
Step 1: Analyze the curve boundary relationships on the given interval.
The curves $y = \sqrt{x}$ and $y = x^2$ intersect at $(0, 0)$ and $(1, 1)$.
On the interval $[\frac{1}{4}, 1]$, we have:
\[
\sqrt{x} \ge x^2
\]
Thus, the upper boundary of the region is $y_{\text{top}} = \sqrt{x}$ and the lower boundary is $y_{\text{bottom}} = x^2$.

Step 2: Set up the integral to find the area.
The area $A$ of the region is:
\[
A = \int_{\frac{1}{4}}^{1} \left( \sqrt{x} - x^2 \right) \, dx = \int_{\frac{1}{4}}^{1} \left( x^{\frac{1}{2}} - x^2 \right) \, dx
\]

Step 3: Integrate and evaluate.
\[
A = \left[ \frac{2}{3} x^{\frac{3}{2}} - \frac{x^3}{3} \right]_{\frac{1}{4}}^{1}
\]
Evaluate at the upper limit $x = 1$:
\[
V(1) = \frac{2}{3}(1) - \frac{1}{3} = \frac{1}{3}
\]
Evaluate at the lower limit $x = \frac{1}{4}$:
\[
V\left(\frac{1}{4}\right) = \frac{2}{3}\left(\frac{1}{4}\right)^{\frac{3}{2}} - \frac{1}{3}\left(\frac{1}{4}\right)^3 = \frac{2}{3}\left(\frac{1}{8}\right) - \frac{1}{192} = \frac{1}{12} - \frac{1}{192} = \frac{16 - 1}{192} = \frac{15}{192} = \frac{5}{64}
\]
Subtract the evaluations:
\[
A = \frac{1}{3} - \frac{5}{64} = \frac{64 - 15}{192} = \frac{49}{192}
\]

Step 4: Sketch description.
The region is bounded above by the curve $y = \sqrt{x}$ (which curves upward from $(1/4, 1/2)$ to $(1,1)$) and below by the parabola $y = x^2$ (which opens upward from $(1/4, 1/16)$ to $(1,1)$), enclosed on the left by the vertical line $x = 1/4$ and closed on the right at their intersection point $(1,1)$.

Final Answer:
\[
\boxed{\frac{49}{192}}
\]

%%%%%%%%%%%%%%%%%%%%%%%%%%%%%%%%%%%%%%%%%%%%%%%%%%
% QUESTION_END
%%%%%%%%%%%%%%%%%%%%%%%%%%%%%%%%%%%%%%%%%%%%%%%%%%


%%%%%%%%%%%%%%%%%%%%%%%%%%%%%%%%%%%%%%%%%%%%%%%%%%
% QUESTION_START
%%%%%%%%%%%%%%%%%%%%%%%%%%%%%%%%%%%%%%%%%%%%%%%%%%

% id=cse23-final-secb-q8b
% discipline=CSE
% year=2023
% exam_type=Term Final
% section=B
% ct_number=UNKNOWN
% question_number=8b
% topic=Definite Integration
% subtopic=General
% difficulty=Hard
% length=Medium
% question_type=Repeated PYQ Pattern
% frequency=1
% appearances=
% OCR_STATUS=VERIFIED

\section*{Term Final - Q8(b)}

Question:
Find the volume of the solid formed by revolving about the x-axis of a loop of the curve $(x - 4a)y^2 = ax(x - 3a)$. [Marks: 05]

Solution:
Step 1: Determine the boundaries of the loop.
Rewrite the curve's equation:
\[
y^2 = \frac{ax(x - 3a)}{x - 4a}
\]
For $y$ to have real values, we must have $y^2 \ge 0$. Assuming $a > 0$, let's analyze the sign of the expression:
\begin{itemize}
    \item When $x < 0 \implies y^2 = \frac{(-)(-)}{(-)} < 0$ (no curve).
    \item When $0 \le x \le 3a \implies y^2 = \frac{(+)(-)}{(-)} \ge 0$ (curve exists, forms a loop).
    \item When $3a < x < 4a \implies y^2 = \frac{(+)(+)}{(-)} < 0$ (no curve).
    \item When $x > 4a \implies y^2 = \frac{(+)(+)}{(+)} > 0$ (curve exists, extends infinitely).
\end{itemize}
Thus, the loop of the curve is defined on the interval $x \in [0, 3a]$.

Step 2: Set up the volume of revolution about the $x$-axis.
The formula for the volume of revolution about the $x$-axis is:
\[
V = \pi \int_{0}^{3a} y^2 \, dx = \pi \int_{0}^{3a} \frac{ax(x-3a)}{x-4a} \, dx = \pi a \int_{0}^{3a} \frac{x^2 - 3ax}{x-4a} \, dx
\]

Step 3: Simplify the integrand using algebraic long division.
Divide $x^2 - 3ax$ by $x - 4a$:
\[
x^2 - 3ax = (x - 4a)(x + a) + 4a^2
\]
Thus:
\[
\frac{x^2 - 3ax}{x - 4a} = x + a + \frac{4a^2}{x - 4a}
\]

Step 4: Integrate and evaluate.
\[
V = \pi a \int_{0}^{3a} \left( x + a + \frac{4a^2}{x - 4a} \right) \, dx
\]
Evaluate the polynomial term integration:
\[
\int_{0}^{3a} (x + a) \, dx = \left[ \frac{x^2}{2} + ax \right]_{0}^{3a} = \frac{9a^2}{2} + 3a^2 = \frac{15a^2}{2}
\]
Evaluate the rational term integration:
\[
\int_{0}^{3a} \frac{4a^2}{x-4a} \, dx = 4a^2 \left[ \log |x - 4a| \right]_{0}^{3a} = 4a^2 \left( \log|a| - \log|4a| \right) = 4a^2 \log\left(\frac{1}{4}\right) = -8a^2 \log 2
\]
Combine both parts:
\[
V = \pi a \left( \frac{15a^2}{2} - 8a^2 \log 2 \right) = \pi a^3 \left( \frac{15}{2} - 8\log 2 \right)
\]

Final Answer:
\[
\boxed{\pi a^3 \left( \frac{15}{2} - 8\log 2 \right)}
\]

%%%%%%%%%%%%%%%%%%%%%%%%%%%%%%%%%%%%%%%%%%%%%%%%%%
% QUESTION_END
%%%%%%%%%%%%%%%%%%%%%%%%%%%%%%%%%%%%%%%%%%%%%%%%%%

\end{document}